\documentclass[]{article}
\usepackage{lmodern}
\usepackage{amssymb,amsmath}
\usepackage{ifxetex,ifluatex}
\usepackage{fixltx2e} % provides \textsubscript
\ifnum 0\ifxetex 1\fi\ifluatex 1\fi=0 % if pdftex
  \usepackage[T1]{fontenc}
  \usepackage[utf8]{inputenc}
\else % if luatex or xelatex
  \ifxetex
    \usepackage{mathspec}
    \usepackage{xltxtra,xunicode}
  \else
    \usepackage{fontspec}
  \fi
  \defaultfontfeatures{Mapping=tex-text,Scale=MatchLowercase}
  \newcommand{\euro}{€}
\fi
% use upquote if available, for straight quotes in verbatim environments
\IfFileExists{upquote.sty}{\usepackage{upquote}}{}
% use microtype if available
\IfFileExists{microtype.sty}{\usepackage{microtype}}{}
\usepackage[margin=1in]{geometry}
\usepackage{color}
\usepackage{fancyvrb}
\newcommand{\VerbBar}{|}
\newcommand{\VERB}{\Verb[commandchars=\\\{\}]}
\DefineVerbatimEnvironment{Highlighting}{Verbatim}{commandchars=\\\{\}}
% Add ',fontsize=\small' for more characters per line
\newenvironment{Shaded}{}{}
\newcommand{\AlertTok}[1]{\textcolor[rgb]{1.00,0.00,0.00}{\textbf{#1}}}
\newcommand{\AnnotationTok}[1]{\textcolor[rgb]{0.38,0.63,0.69}{\textbf{\textit{#1}}}}
\newcommand{\AttributeTok}[1]{\textcolor[rgb]{0.49,0.56,0.16}{#1}}
\newcommand{\BaseNTok}[1]{\textcolor[rgb]{0.25,0.63,0.44}{#1}}
\newcommand{\BuiltInTok}[1]{\textcolor[rgb]{0.00,0.50,0.00}{#1}}
\newcommand{\CharTok}[1]{\textcolor[rgb]{0.25,0.44,0.63}{#1}}
\newcommand{\CommentTok}[1]{\textcolor[rgb]{0.38,0.63,0.69}{\textit{#1}}}
\newcommand{\CommentVarTok}[1]{\textcolor[rgb]{0.38,0.63,0.69}{\textbf{\textit{#1}}}}
\newcommand{\ConstantTok}[1]{\textcolor[rgb]{0.53,0.00,0.00}{#1}}
\newcommand{\ControlFlowTok}[1]{\textcolor[rgb]{0.00,0.44,0.13}{\textbf{#1}}}
\newcommand{\DataTypeTok}[1]{\textcolor[rgb]{0.56,0.13,0.00}{#1}}
\newcommand{\DecValTok}[1]{\textcolor[rgb]{0.25,0.63,0.44}{#1}}
\newcommand{\DocumentationTok}[1]{\textcolor[rgb]{0.73,0.13,0.13}{\textit{#1}}}
\newcommand{\ErrorTok}[1]{\textcolor[rgb]{1.00,0.00,0.00}{\textbf{#1}}}
\newcommand{\ExtensionTok}[1]{#1}
\newcommand{\FloatTok}[1]{\textcolor[rgb]{0.25,0.63,0.44}{#1}}
\newcommand{\FunctionTok}[1]{\textcolor[rgb]{0.02,0.16,0.49}{#1}}
\newcommand{\ImportTok}[1]{\textcolor[rgb]{0.00,0.50,0.00}{\textbf{#1}}}
\newcommand{\InformationTok}[1]{\textcolor[rgb]{0.38,0.63,0.69}{\textbf{\textit{#1}}}}
\newcommand{\KeywordTok}[1]{\textcolor[rgb]{0.00,0.44,0.13}{\textbf{#1}}}
\newcommand{\NormalTok}[1]{#1}
\newcommand{\OperatorTok}[1]{\textcolor[rgb]{0.40,0.40,0.40}{#1}}
\newcommand{\OtherTok}[1]{\textcolor[rgb]{0.00,0.44,0.13}{#1}}
\newcommand{\PreprocessorTok}[1]{\textcolor[rgb]{0.74,0.48,0.00}{#1}}
\newcommand{\RegionMarkerTok}[1]{#1}
\newcommand{\SpecialCharTok}[1]{\textcolor[rgb]{0.25,0.44,0.63}{#1}}
\newcommand{\SpecialStringTok}[1]{\textcolor[rgb]{0.73,0.40,0.53}{#1}}
\newcommand{\StringTok}[1]{\textcolor[rgb]{0.25,0.44,0.63}{#1}}
\newcommand{\VariableTok}[1]{\textcolor[rgb]{0.10,0.09,0.49}{#1}}
\newcommand{\VerbatimStringTok}[1]{\textcolor[rgb]{0.25,0.44,0.63}{#1}}
\newcommand{\WarningTok}[1]{\textcolor[rgb]{0.38,0.63,0.69}{\textbf{\textit{#1}}}}
\ifxetex
  \usepackage[setpagesize=false, % page size defined by xetex
              unicode=false, % unicode breaks when used with xetex
              xetex]{hyperref}
\else
  \usepackage[unicode=true]{hyperref}
\fi
\hypersetup{breaklinks=true,
            bookmarks=true,
            pdfauthor={},
            pdftitle={Faking ADTs and GADTs in Languages That Shouldn't Have Them},
            colorlinks=true,
            citecolor=blue,
            urlcolor=blue,
            linkcolor=magenta,
            pdfborder={0 0 0}}
\urlstyle{same}  % don't use monospace font for urls
% Make links footnotes instead of hotlinks:
\renewcommand{\href}[2]{#2\footnote{\url{#1}}}
\setlength{\parindent}{0pt}
\setlength{\parskip}{6pt plus 2pt minus 1pt}
\setlength{\emergencystretch}{3em}  % prevent overfull lines
\setcounter{secnumdepth}{0}

\title{Faking ADTs and GADTs in Languages That Shouldn't Have Them}

\begin{document}
\maketitle

\% Justin Le

\emph{Originally posted on
\textbf{\href{https://blog.jle.im/entry/faking-adts-gadts-in-languages-that-shouldnt-have-them.html}{in
Code}}.}

Haskell is the world's best programming language\footnote{I bet you thought
  there was going be some sort of caveat in this footnote, didn't you?
  {[}kmett{]}: https://x.com/kmett/status/1844812186608099463}, but let's face
the harsh reality that a lot of times in life you'll have to write in other
programming languages. But alas you have been fully
{[}Haskell-brained{]}{[}kmett{]} and lost all ability to program unless it is
type-directed, you don't even know how to start writing a program without
imagining its shape as a type first.

Well, fear not. The foundational theory behind ADTs and GADTs are so fundamental
that they'll fit (somewhat) seamlessly into whatever language you're forced to
write. After all, if they can fit
\href{https://www.reddit.com/r/haskell/comments/9m2o5r/digging_reveals_profunctor_optics_in_mineacraft/}{profunctor
optics in Microsoft's Java code}, the sky's the limit!

\section{ADTs and the Visitor Pattern}\label{adts-and-the-visitor-pattern}

Let's get normal ADT's out of the way. Most languages do have \emph{structs}.
Not all languages have \emph{immutable} structs, but we'll take what we can get.
With immutable structs we get product types.

Examples:

\begin{enumerate}
\def\labelenumi{\arabic{enumi}.}
\tightlist
\item
  x,y tuples
\item
  one-way message protocol, or state machine? tree? untyped expr?
\item
  lists
\item
  vector
\item
  expr
\item
  message protocol
\end{enumerate}

\subsection{Implementation}\label{implementation}

Typically if your language does not natively support sum types, you can achieve
the same design patterns and guarantees using the ``visitor pattern'', which is
essentially the ``case match'' as a data structure or argument list.

If you don't have subtyping and inheritance, usually this looks like defining
some sort of union type with a tag.

I'm going to use C++ \emph{without classes} or the stdlib (using only raw
function pointers and structs) to demonstrate how you'd do this with generic
types:

\begin{Shaded}
\begin{Highlighting}[]
\KeywordTok{template} \OperatorTok{\textless{}}\KeywordTok{typename}\NormalTok{ T}\OperatorTok{\textgreater{}}
\KeywordTok{struct}\NormalTok{ Maybe }\OperatorTok{\{}
    \DataTypeTok{bool}\NormalTok{ present}\OperatorTok{;}
\NormalTok{    T value}\OperatorTok{;}
\OperatorTok{\};}

\KeywordTok{template} \OperatorTok{\textless{}}\KeywordTok{typename}\NormalTok{ T}\OperatorTok{,} \KeywordTok{typename}\NormalTok{ R}\OperatorTok{\textgreater{}}
\KeywordTok{struct}\NormalTok{ MaybeVisitor }\OperatorTok{\{}
\NormalTok{    R }\OperatorTok{(*}\NormalTok{visitJust}\OperatorTok{)(}\NormalTok{T}\OperatorTok{);}
\NormalTok{    R }\OperatorTok{(*}\NormalTok{visitNothing}\OperatorTok{)();}
\OperatorTok{\};}

\KeywordTok{template} \OperatorTok{\textless{}}\KeywordTok{typename}\NormalTok{ T}\OperatorTok{,} \KeywordTok{typename}\NormalTok{ R}\OperatorTok{\textgreater{}}
\NormalTok{R acceptMaybe}\OperatorTok{(}\NormalTok{Maybe}\OperatorTok{\textless{}}\NormalTok{T}\OperatorTok{\textgreater{}}\NormalTok{ maybe}\OperatorTok{,}\NormalTok{ MaybeVisitor}\OperatorTok{\textless{}}\NormalTok{T}\OperatorTok{,}\NormalTok{ R}\OperatorTok{\textgreater{}}\NormalTok{ visitor}\OperatorTok{)} \OperatorTok{\{}
    \ControlFlowTok{if} \OperatorTok{(}\NormalTok{maybe}\OperatorTok{.}\NormalTok{present}\OperatorTok{)} \OperatorTok{\{}
        \ControlFlowTok{return}\NormalTok{ visitor}\OperatorTok{.}\NormalTok{visitJust}\OperatorTok{(}\NormalTok{maybe}\OperatorTok{.}\NormalTok{value}\OperatorTok{);}
    \OperatorTok{\}} \ControlFlowTok{else} \OperatorTok{\{}
        \ControlFlowTok{return}\NormalTok{ visitor}\OperatorTok{.}\NormalTok{visitNothing}\OperatorTok{();}
    \OperatorTok{\}}
\OperatorTok{\}}
\end{Highlighting}
\end{Shaded}

If you have multiple types of payloads (instead of just \texttt{T}) then you can
use a union to save space. This is why sum types are often called `tagged
unions': they are a tag with a union!

You can create the values using:

\begin{Shaded}
\begin{Highlighting}[]
\KeywordTok{template} \OperatorTok{\textless{}}\KeywordTok{typename}\NormalTok{ T}\OperatorTok{\textgreater{}}
\NormalTok{Maybe}\OperatorTok{\textless{}}\NormalTok{T}\OperatorTok{\textgreater{}}\NormalTok{ mkJust}\OperatorTok{(}\NormalTok{T value}\OperatorTok{)} \OperatorTok{\{} \ControlFlowTok{return} \OperatorTok{\{} \KeywordTok{true}\OperatorTok{,}\NormalTok{ value }\OperatorTok{\};} \OperatorTok{\}}

\KeywordTok{template} \OperatorTok{\textless{}}\KeywordTok{typename}\NormalTok{ T}\OperatorTok{\textgreater{}}
\NormalTok{Maybe}\OperatorTok{\textless{}}\NormalTok{T}\OperatorTok{\textgreater{}}\NormalTok{ mkNothing}\OperatorTok{()} \OperatorTok{\{} \ControlFlowTok{return} \OperatorTok{\{} \KeywordTok{false} \OperatorTok{\};} \OperatorTok{\}}
\end{Highlighting}
\end{Shaded}

And, to implement \texttt{showMaybeInt}, we'd create a
\texttt{MaybeVisitor\textless{}int,std::string\textgreater{}}:

\begin{Shaded}
\begin{Highlighting}[]
\BuiltInTok{std::}\NormalTok{string showMaybeInt}\OperatorTok{(}\NormalTok{Maybe}\OperatorTok{\textless{}}\DataTypeTok{int}\OperatorTok{\textgreater{}}\NormalTok{ maybe}\OperatorTok{)} \OperatorTok{\{}
    \ControlFlowTok{return}\NormalTok{ acceptMaybe}\OperatorTok{(}\NormalTok{maybe}\OperatorTok{,} \OperatorTok{\{}
        \OperatorTok{[](}\DataTypeTok{int}\NormalTok{ value}\OperatorTok{)} \OperatorTok{\{} \ControlFlowTok{return} \StringTok{"Something is here: "} \OperatorTok{+} \BuiltInTok{std::}\NormalTok{to\_string}\OperatorTok{(}\NormalTok{value}\OperatorTok{);} \OperatorTok{\},}
        \OperatorTok{[]()} \OperatorTok{\{} \ControlFlowTok{return} \StringTok{"There\textquotesingle{}s nothing here"}\OperatorTok{;} \OperatorTok{\}}
    \OperatorTok{\});}
\OperatorTok{\}}
\end{Highlighting}
\end{Shaded}

Note that in this way, the compiler enforces that we handle every branch. And,
if we ever add a new branch, everything that ever consumes \texttt{MaybeVisitor}
will have to add a new handler.

TODO: use Command to illustrate the union

However, if your language as subtyping built in (maybe with classes and
subclasses), you can implement it in terms of that, which is nice in python,
java, etc.

\begin{Shaded}
\begin{Highlighting}[]
\KeywordTok{interface}\NormalTok{ ExprVisitor}\OperatorTok{\textless{}}\NormalTok{R}\OperatorTok{\textgreater{}} \OperatorTok{\{}
\NormalTok{    R }\FunctionTok{visitLit}\OperatorTok{(}\DataTypeTok{int}\NormalTok{ value}\OperatorTok{);}
\NormalTok{    R }\FunctionTok{visitNegate}\OperatorTok{(}\NormalTok{Expr expr}\OperatorTok{);}
\NormalTok{    R }\FunctionTok{visitAdd}\OperatorTok{(}\NormalTok{Expr left}\OperatorTok{,}\NormalTok{ Expr right}\OperatorTok{);}
\NormalTok{    R }\FunctionTok{visitSub}\OperatorTok{(}\NormalTok{Expr left}\OperatorTok{,}\NormalTok{ Expr right}\OperatorTok{);}
\NormalTok{    R }\FunctionTok{visitMul}\OperatorTok{(}\NormalTok{Expr left}\OperatorTok{,}\NormalTok{ Expr right}\OperatorTok{);}
\OperatorTok{\}}

\KeywordTok{abstract} \KeywordTok{class}\NormalTok{ Expr }\OperatorTok{\{}
    \KeywordTok{public} \KeywordTok{abstract} \OperatorTok{\textless{}}\NormalTok{R}\OperatorTok{\textgreater{}}\NormalTok{ R }\FunctionTok{accept}\OperatorTok{(}\NormalTok{ExprVisitor}\OperatorTok{\textless{}}\NormalTok{R}\OperatorTok{\textgreater{}}\NormalTok{ visitor}\OperatorTok{);}
\OperatorTok{\}}
\end{Highlighting}
\end{Shaded}

Now instead of manually implementing a tag, we take advantage of our language's
default dynamic dispatch to handle it for us. Each constructor becomes a class,
but it's important to \emph{only allow} access using \texttt{accept} to properly
enforce the sum type pattern.

\begin{Shaded}
\begin{Highlighting}[]
\KeywordTok{class}\NormalTok{ Lit }\KeywordTok{extends}\NormalTok{ Expr }\OperatorTok{\{}
    \KeywordTok{private} \DataTypeTok{final} \DataTypeTok{int}\NormalTok{ value}\OperatorTok{;}

    \KeywordTok{public} \FunctionTok{Lit}\OperatorTok{(}\DataTypeTok{int}\NormalTok{ value}\OperatorTok{)} \OperatorTok{\{}
        \KeywordTok{this}\OperatorTok{.}\FunctionTok{value} \OperatorTok{=}\NormalTok{ value}\OperatorTok{;}
    \OperatorTok{\}}

    \AttributeTok{@Override}
    \KeywordTok{public} \OperatorTok{\textless{}}\NormalTok{R}\OperatorTok{\textgreater{}}\NormalTok{ R }\FunctionTok{accept}\OperatorTok{(}\NormalTok{ExprVisitor}\OperatorTok{\textless{}}\NormalTok{R}\OperatorTok{\textgreater{}}\NormalTok{ visitor}\OperatorTok{)} \OperatorTok{\{}
        \ControlFlowTok{return}\NormalTok{ visitor}\OperatorTok{.}\FunctionTok{visitLit}\OperatorTok{(}\NormalTok{value}\OperatorTok{);}
    \OperatorTok{\}}
\OperatorTok{\}}

\KeywordTok{class}\NormalTok{ Negate }\KeywordTok{extends}\NormalTok{ Expr }\OperatorTok{\{}
    \KeywordTok{private} \DataTypeTok{final}\NormalTok{ Expr expr}\OperatorTok{;}

    \KeywordTok{public} \FunctionTok{Negate}\OperatorTok{(}\NormalTok{Expr expr}\OperatorTok{)} \OperatorTok{\{} \KeywordTok{this}\OperatorTok{.}\FunctionTok{expr} \OperatorTok{=}\NormalTok{ expr}\OperatorTok{;} \OperatorTok{\}}

    \AttributeTok{@Override}
    \KeywordTok{public} \OperatorTok{\textless{}}\NormalTok{R}\OperatorTok{\textgreater{}}\NormalTok{ R }\FunctionTok{accept}\OperatorTok{(}\NormalTok{ExprVisitor}\OperatorTok{\textless{}}\NormalTok{R}\OperatorTok{\textgreater{}}\NormalTok{ visitor}\OperatorTok{)} \OperatorTok{\{}
        \ControlFlowTok{return}\NormalTok{ visitor}\OperatorTok{.}\FunctionTok{visitNegate}\OperatorTok{(}\NormalTok{expr}\OperatorTok{);}
    \OperatorTok{\}}
\OperatorTok{\}}

\KeywordTok{class}\NormalTok{ Add }\KeywordTok{extends}\NormalTok{ Expr }\OperatorTok{\{}
    \KeywordTok{private} \DataTypeTok{final}\NormalTok{ Expr left}\OperatorTok{;}
    \KeywordTok{private} \DataTypeTok{final}\NormalTok{ Expr right}\OperatorTok{;}

    \KeywordTok{public} \FunctionTok{Add}\OperatorTok{(}\NormalTok{Expr left}\OperatorTok{,}\NormalTok{ Expr right}\OperatorTok{)} \OperatorTok{\{}
        \KeywordTok{this}\OperatorTok{.}\FunctionTok{left} \OperatorTok{=}\NormalTok{ left}\OperatorTok{;}
        \KeywordTok{this}\OperatorTok{.}\FunctionTok{right} \OperatorTok{=}\NormalTok{ right}\OperatorTok{;}
    \OperatorTok{\}}

    \AttributeTok{@Override}
    \KeywordTok{public} \OperatorTok{\textless{}}\NormalTok{R}\OperatorTok{\textgreater{}}\NormalTok{ R }\FunctionTok{accept}\OperatorTok{(}\NormalTok{ExprVisitor}\OperatorTok{\textless{}}\NormalTok{R}\OperatorTok{\textgreater{}}\NormalTok{ visitor}\OperatorTok{)} \OperatorTok{\{}
        \ControlFlowTok{return}\NormalTok{ visitor}\OperatorTok{.}\FunctionTok{visitAdd}\OperatorTok{(}\NormalTok{left}\OperatorTok{,}\NormalTok{ right}\OperatorTok{);}
    \OperatorTok{\}}
\OperatorTok{\}}

\KeywordTok{class}\NormalTok{ Sub }\KeywordTok{extends}\NormalTok{ Expr }\OperatorTok{\{}
    \KeywordTok{private} \DataTypeTok{final}\NormalTok{ Expr left}\OperatorTok{;}
    \KeywordTok{private} \DataTypeTok{final}\NormalTok{ Expr right}\OperatorTok{;}

    \KeywordTok{public} \FunctionTok{Sub}\OperatorTok{(}\NormalTok{Expr left}\OperatorTok{,}\NormalTok{ Expr right}\OperatorTok{)} \OperatorTok{\{}
        \KeywordTok{this}\OperatorTok{.}\FunctionTok{left} \OperatorTok{=}\NormalTok{ left}\OperatorTok{;}
        \KeywordTok{this}\OperatorTok{.}\FunctionTok{right} \OperatorTok{=}\NormalTok{ right}\OperatorTok{;}
    \OperatorTok{\}}

    \AttributeTok{@Override}
    \KeywordTok{public} \OperatorTok{\textless{}}\NormalTok{R}\OperatorTok{\textgreater{}}\NormalTok{ R }\FunctionTok{accept}\OperatorTok{(}\NormalTok{ExprVisitor}\OperatorTok{\textless{}}\NormalTok{R}\OperatorTok{\textgreater{}}\NormalTok{ visitor}\OperatorTok{)} \OperatorTok{\{}
        \ControlFlowTok{return}\NormalTok{ visitor}\OperatorTok{.}\FunctionTok{visitSub}\OperatorTok{(}\NormalTok{left}\OperatorTok{,}\NormalTok{ right}\OperatorTok{);}
    \OperatorTok{\}}
\OperatorTok{\}}

\KeywordTok{class}\NormalTok{ Mul }\KeywordTok{extends}\NormalTok{ Expr }\OperatorTok{\{}
    \KeywordTok{private} \DataTypeTok{final}\NormalTok{ Expr left}\OperatorTok{;}
    \KeywordTok{private} \DataTypeTok{final}\NormalTok{ Expr right}\OperatorTok{;}

    \KeywordTok{public} \FunctionTok{Mul}\OperatorTok{(}\NormalTok{Expr left}\OperatorTok{,}\NormalTok{ Expr right}\OperatorTok{)} \OperatorTok{\{}
        \KeywordTok{this}\OperatorTok{.}\FunctionTok{left} \OperatorTok{=}\NormalTok{ left}\OperatorTok{;}
        \KeywordTok{this}\OperatorTok{.}\FunctionTok{right} \OperatorTok{=}\NormalTok{ right}\OperatorTok{;}
    \OperatorTok{\}}

    \AttributeTok{@Override}
    \KeywordTok{public} \OperatorTok{\textless{}}\NormalTok{R}\OperatorTok{\textgreater{}}\NormalTok{ R }\FunctionTok{accept}\OperatorTok{(}\NormalTok{ExprVisitor}\OperatorTok{\textless{}}\NormalTok{R}\OperatorTok{\textgreater{}}\NormalTok{ visitor}\OperatorTok{)} \OperatorTok{\{}
        \ControlFlowTok{return}\NormalTok{ visitor}\OperatorTok{.}\FunctionTok{visitMul}\OperatorTok{(}\NormalTok{left}\OperatorTok{,}\NormalTok{ right}\OperatorTok{);}
    \OperatorTok{\}}
\OperatorTok{\}}
\end{Highlighting}
\end{Shaded}

Alternatively you could make all of the subclasses anonymous and expose them as
factory methods, if your language allows it:

\begin{Shaded}
\begin{Highlighting}[]
\KeywordTok{abstract} \KeywordTok{class}\NormalTok{ Expr }\OperatorTok{\{}
    \KeywordTok{public} \KeywordTok{abstract} \OperatorTok{\textless{}}\NormalTok{R}\OperatorTok{\textgreater{}}\NormalTok{ R }\FunctionTok{accept}\OperatorTok{(}\NormalTok{ExprVisitor}\OperatorTok{\textless{}}\NormalTok{R}\OperatorTok{\textgreater{}}\NormalTok{ visitor}\OperatorTok{);}

    \KeywordTok{public} \DataTypeTok{static}\NormalTok{ Expr }\FunctionTok{lit}\OperatorTok{(}\DataTypeTok{int}\NormalTok{ value}\OperatorTok{)} \OperatorTok{\{}
        \ControlFlowTok{return} \KeywordTok{new} \FunctionTok{Expr}\OperatorTok{()} \OperatorTok{\{}
            \AttributeTok{@Override}
            \KeywordTok{public} \OperatorTok{\textless{}}\NormalTok{R}\OperatorTok{\textgreater{}}\NormalTok{ R }\FunctionTok{accept}\OperatorTok{(}\NormalTok{ExprVisitor}\OperatorTok{\textless{}}\NormalTok{R}\OperatorTok{\textgreater{}}\NormalTok{ visitor}\OperatorTok{)} \OperatorTok{\{}
                \ControlFlowTok{return}\NormalTok{ visitor}\OperatorTok{.}\FunctionTok{visitLit}\OperatorTok{(}\NormalTok{value}\OperatorTok{);}
            \OperatorTok{\}}
        \OperatorTok{\};}
    \OperatorTok{\}}

    \KeywordTok{public} \DataTypeTok{static}\NormalTok{ Expr }\FunctionTok{negate}\OperatorTok{(}\NormalTok{Expr unary}\OperatorTok{)} \OperatorTok{\{}
        \ControlFlowTok{return} \KeywordTok{new} \FunctionTok{Expr}\OperatorTok{()} \OperatorTok{\{}
            \AttributeTok{@Override}
            \KeywordTok{public} \OperatorTok{\textless{}}\NormalTok{R}\OperatorTok{\textgreater{}}\NormalTok{ R }\FunctionTok{accept}\OperatorTok{(}\NormalTok{ExprVisitor}\OperatorTok{\textless{}}\NormalTok{R}\OperatorTok{\textgreater{}}\NormalTok{ visitor}\OperatorTok{)} \OperatorTok{\{}
                \ControlFlowTok{return}\NormalTok{ visitor}\OperatorTok{.}\FunctionTok{visitNegate}\OperatorTok{(}\NormalTok{unary}\OperatorTok{);}
            \OperatorTok{\}}
        \OperatorTok{\};}
    \OperatorTok{\}}

    \KeywordTok{public} \DataTypeTok{static}\NormalTok{ Expr }\FunctionTok{add}\OperatorTok{(}\NormalTok{Expr left}\OperatorTok{,}\NormalTok{ Expr right}\OperatorTok{)} \OperatorTok{\{}
        \ControlFlowTok{return} \KeywordTok{new} \FunctionTok{Expr}\OperatorTok{()} \OperatorTok{\{}
            \AttributeTok{@Override}
            \KeywordTok{public} \OperatorTok{\textless{}}\NormalTok{R}\OperatorTok{\textgreater{}}\NormalTok{ R }\FunctionTok{accept}\OperatorTok{(}\NormalTok{ExprVisitor}\OperatorTok{\textless{}}\NormalTok{R}\OperatorTok{\textgreater{}}\NormalTok{ visitor}\OperatorTok{)} \OperatorTok{\{}
                \ControlFlowTok{return}\NormalTok{ visitor}\OperatorTok{.}\FunctionTok{visitAdd}\OperatorTok{(}\NormalTok{left}\OperatorTok{,}\NormalTok{ right}\OperatorTok{);}
            \OperatorTok{\}}
        \OperatorTok{\};}
    \OperatorTok{\}}

    \CommentTok{// ... etc}
\OperatorTok{\}}
\end{Highlighting}
\end{Shaded}

You'd then call using:

\begin{Shaded}
\begin{Highlighting}[]
\KeywordTok{public} \KeywordTok{class}\NormalTok{ Main }\OperatorTok{\{}
    \KeywordTok{public} \DataTypeTok{static} \DataTypeTok{void} \FunctionTok{main}\OperatorTok{(}\BuiltInTok{String}\OperatorTok{[]}\NormalTok{ args}\OperatorTok{)} \OperatorTok{\{}
\NormalTok{        Expr expr }\OperatorTok{=} \KeywordTok{new} \FunctionTok{Mul}\OperatorTok{(}\KeywordTok{new} \FunctionTok{Negate}\OperatorTok{(}\KeywordTok{new} \FunctionTok{Add}\OperatorTok{(}\KeywordTok{new} \FunctionTok{Lit}\OperatorTok{(}\DecValTok{4}\OperatorTok{),} \KeywordTok{new} \FunctionTok{Lit}\OperatorTok{(}\DecValTok{5}\OperatorTok{))),} \KeywordTok{new} \FunctionTok{Lit}\OperatorTok{(}\DecValTok{8}\OperatorTok{));}
        \CommentTok{// or}
        \CommentTok{// Expr expr = Eval.mul(Eval.negate(Eval.add(Eval.lit(4), Eval.lit(5))), Eval.lit(8));}

\NormalTok{        ExprVisitor}\OperatorTok{\textless{}}\BuiltInTok{Integer}\OperatorTok{\textgreater{}}\NormalTok{ eval }\OperatorTok{=} \KeywordTok{new}\NormalTok{ ExprVisitor}\OperatorTok{\textless{}\textgreater{}()} \OperatorTok{\{}
            \AttributeTok{@Override} \KeywordTok{public} \BuiltInTok{Integer} \FunctionTok{visitLit}\OperatorTok{(}\DataTypeTok{int}\NormalTok{ value}\OperatorTok{)} \OperatorTok{\{} \ControlFlowTok{return}\NormalTok{ value}\OperatorTok{;} \OperatorTok{\}}
            \AttributeTok{@Override} \KeywordTok{public} \BuiltInTok{Integer} \FunctionTok{visitNegate}\OperatorTok{(}\NormalTok{Expr expr}\OperatorTok{)} \OperatorTok{\{} \ControlFlowTok{return} \OperatorTok{{-}}\NormalTok{expr}\OperatorTok{.}\FunctionTok{accept}\OperatorTok{(}\KeywordTok{this}\OperatorTok{);} \OperatorTok{\}}
            \AttributeTok{@Override} \KeywordTok{public} \BuiltInTok{Integer} \FunctionTok{visitAdd}\OperatorTok{(}\NormalTok{Expr left}\OperatorTok{,}\NormalTok{ Expr right}\OperatorTok{)} \OperatorTok{\{} \ControlFlowTok{return}\NormalTok{ left}\OperatorTok{.}\FunctionTok{accept}\OperatorTok{(}\KeywordTok{this}\OperatorTok{)} \OperatorTok{+}\NormalTok{ right}\OperatorTok{.}\FunctionTok{accept}\OperatorTok{(}\KeywordTok{this}\OperatorTok{);} \OperatorTok{\}}
            \AttributeTok{@Override} \KeywordTok{public} \BuiltInTok{Integer} \FunctionTok{visitSub}\OperatorTok{(}\NormalTok{Expr left}\OperatorTok{,}\NormalTok{ Expr right}\OperatorTok{)} \OperatorTok{\{} \ControlFlowTok{return}\NormalTok{ left}\OperatorTok{.}\FunctionTok{accept}\OperatorTok{(}\KeywordTok{this}\OperatorTok{)} \OperatorTok{{-}}\NormalTok{ right}\OperatorTok{.}\FunctionTok{accept}\OperatorTok{(}\KeywordTok{this}\OperatorTok{);} \OperatorTok{\}}
            \AttributeTok{@Override} \KeywordTok{public} \BuiltInTok{Integer} \FunctionTok{visitMul}\OperatorTok{(}\NormalTok{Expr left}\OperatorTok{,}\NormalTok{ Expr right}\OperatorTok{)} \OperatorTok{\{} \ControlFlowTok{return}\NormalTok{ left}\OperatorTok{.}\FunctionTok{accept}\OperatorTok{(}\KeywordTok{this}\OperatorTok{)} \OperatorTok{*}\NormalTok{ right}\OperatorTok{.}\FunctionTok{accept}\OperatorTok{(}\KeywordTok{this}\OperatorTok{);} \OperatorTok{\}}
        \OperatorTok{\};}

        \BuiltInTok{System}\OperatorTok{.}\FunctionTok{out}\OperatorTok{.}\FunctionTok{println}\OperatorTok{(}\StringTok{"Result: "} \OperatorTok{+}\NormalTok{ expr}\OperatorTok{.}\FunctionTok{accept}\OperatorTok{(}\NormalTok{eval}\OperatorTok{));}
    \OperatorTok{\}}
\OperatorTok{\}}
\end{Highlighting}
\end{Shaded}

\section{Signoff}\label{signoff}

Hi, thanks for reading! You can reach me via email at
\href{mailto:justin@jle.im}{\nolinkurl{justin@jle.im}}, or at twitter at
\href{https://twitter.com/mstk}{@mstk}! This post and all others are published
under the \href{https://creativecommons.org/licenses/by-nc-nd/3.0/}{CC-BY-NC-ND
3.0} license. Corrections and edits via pull request are welcome and encouraged
at \href{https://github.com/mstksg/inCode}{the source repository}.

If you feel inclined, or this post was particularly helpful for you, why not
consider \href{https://www.patreon.com/justinle/overview}{supporting me on
Patreon}, or a \href{bitcoin:3D7rmAYgbDnp4gp4rf22THsGt74fNucPDU}{BTC donation}?
:)

\end{document}
